\documentclass[11pt]{article}

\usepackage[utf8]{inputenc}
\usepackage[T1]{fontenc}
\usepackage[spanish,es-noshorthands,es-tabla]{babel}
\usepackage[a4paper,margin=2.05cm]{geometry}
\usepackage{amsmath,amssymb}
\usepackage{booktabs}
\usepackage{longtable}
\usepackage{tabularx}
\usepackage{array}
\usepackage{xcolor}
\usepackage{hyperref}

\hypersetup{
  colorlinks=true,
  linkcolor=blue!55!black,
  urlcolor=blue!55!black,
  citecolor=blue!55!black
}

\setlength{\emergencystretch}{2em}
\renewcommand{\arraystretch}{1.12}

\newcolumntype{L}[1]{>{\raggedright\arraybackslash}p{#1}}
\newcommand{\codepath}[1]{\begingroup\Urlmuskip=0mu plus 2mu\path{#1}\endgroup}
\newcommand{\status}[1]{\texttt{\small #1}}
\newcommand{\target}{\texttt{TARGET}}
\newcommand{\hiddenverified}{\texttt{hidden\_verified}}
\newcommand{\qmain}{0.9998}
\newcommand{\Ezero}{E_0}
\newcommand{\Eplus}{E_+}
\newcommand{\Eminus}{E_-}

\title{Resumen de verificaciones de ocultedad y candidatos Lur'e/Machado}
\author{Documento generado desde salidas CSV/JSON del repositorio}
\date{Datos consultados al 18 de mayo de 2026; salidas principales del 17 de mayo de 2026}

\begin{document}
\maketitle

\begin{abstract}
Este documento resume los datos disponibles hasta ahora para tres bloques:
la replicacion/verificacion del atractor reportado por Danca en el sistema de
Chua fraccionario no suave, los mejores candidatos generados por la ruta
Lur'e/funcion descriptiva sesgada, y los mejores candidatos generados por la
extension tipo Machado. No se incluyen imagenes porque las figuras aun estan
en ajuste. La lectura central es deliberadamente conservadora: ninguna salida
actual debe llamarse \hiddenverified. Las semillas armonicas produjeron
atractores acotados no triviales, pero las pruebas de vecindad de equilibrio
detectaron contactos \target{} en casos clave.
\end{abstract}

\section{Convencion de lectura}

En este proyecto se separan cuatro niveles:
\[
\text{prediccion armonica}
\;\Longrightarrow\;
\text{semilla numerica}
\;\Longrightarrow\;
\text{atractor observado en Caputo}
\;\Longrightarrow\;
\text{verificacion de ocultedad}.
\]
La funcion descriptiva clasica o tipo Machado solo ocupa los dos primeros
niveles. Una trayectoria clasificada como \target{} desde una vecindad pequena
de algun equilibrio bloquea la interpretacion como atractor oculto bajo el
contrato probado. Una salida sin contactos \target{} se reporta como
compatibilidad bajo la muestra y radios evaluados, no como demostracion.

\begin{table}[h]
\centering
\caption{Etiquetas usadas en las salidas y su interpretacion.}
\small
\begin{tabularx}{\textwidth}{L{0.29\textwidth}L{0.64\textwidth}}
\toprule
Etiqueta & Interpretacion cientifica \\
\midrule
\status{target\_attractor} o \target{} &
La trayectoria fue clasificada como perteneciente al atractor candidato o a su
geometria de referencia. Si parte de una vecindad de equilibrio, es evidencia
en contra de ocultedad. \\
\status{compatible\_with\_hiddenness\_under\_tested\_radii} &
No se encontraron contactos \target{} en las trayectorias completadas. Es
evidencia operativa limitada; no equivale a \hiddenverified. \\
\status{compatible\_with\_hiddenness\_under\_tested\_spheres} &
No hubo contactos \target{} en cascarones esfericos discretos. No descarta
contactos dirigidos o no muestreados. \\
\status{not\_supported\_by\_...} &
Se encontro por lo menos una trayectoria desde vecindad de equilibrio que llega
al candidato. La ocultedad no queda soportada por esa prueba. \\
\status{unknown} &
La trayectoria no se pudo clasificar con el contrato actual; requiere mayor
tiempo, refinamiento o criterio geometrico adicional. \\
\bottomrule
\end{tabularx}
\end{table}

\section{Resumen directo de condiciones iniciales y distancias}

La tabla siguiente concentra, antes del detalle metodologico, cuantas
condiciones iniciales o trayectorias se lanzaron, cuantas salieron \target{} y
que tan cerca estaban de los equilibrios al iniciar. Cuando la prueba no parte
de vecindades de equilibrio, se marca como ``no aplica''.

\begin{longtable}{@{}L{0.25\textwidth}L{0.13\textwidth}L{0.11\textwidth}L{0.22\textwidth}L{0.23\textwidth}@{}}
\caption{Conteo operativo de condiciones iniciales, targets y distancia al equilibrio.}\\
\toprule
Prueba & Trayectorias/CI & \target{} & Distancia inicial a equilibrio & Lectura \\
\midrule
\endfirsthead
\toprule
Prueba & Trayectorias/CI & \target{} & Distancia inicial a equilibrio & Lectura \\
\midrule
\endhead
Replica local Danca ABM &
200 & 102 &
\(0.01\) desde \(\Eplus\) y \(\Eminus\) &
Falla la prueba de ocultedad local: 54/100 targets desde \(\Eplus\) y 48/100
desde \(\Eminus\). \\
Corte de cuenca Danca ABM &
10201 & 976 &
No es una prueba local; corte \(xy\) completo &
Sirve para ver cuenca, no para decidir ocultedad por si solo. \\
Corte de cuenca proyecto EFORK &
10201 & 3 &
No es una prueba local; corte \(xy\) completo &
Clasificador con memoria finita; muchos puntos terminan en infinito o
desconocido. \\
Filtro temprano Lur'e &
25 direcciones unicas, 28 filas crudas & 13 direcciones &
\(10^{-5}\) desde \(\Eplus,\Eminus,\Ezero\) &
Bloquea \hiddenverified{} para \status{rank\_0001} y \status{rank\_0004}. \\
Corrida 1 refinada Machado &
1622 & 0 &
Bola/cascaron en \(\Ezero\), radios \(10^{-5}\) y \(3\times10^{-5}\) &
Compatible solo bajo los radios completados; no evalua completamente el segundo
candidato. \\
Machado dirigido LM10 &
2 dirigidas + 2 reproducciones & 2 dirigidas + 2 reproducidas &
\(10^{-5}\) desde \(\Eminus\), direccion \status{axis\_p\_x} &
Bloquea ocultedad para los dos candidatos Machado top. \\
Top-3 cascarones esfericos &
4500 & 0 &
Radios \(10^{-5},3\times10^{-5},10^{-4},3\times10^{-4},10^{-3}\) desde
\(\Ezero,\Eplus,\Eminus\) &
No observa target en cascarones aleatorios, pero no cancela los contactos
dirigidos. \\
Robustez por reclasificacion top-3 &
12 & 0 &
No aplica; parte de estados del candidato &
No robusto como target bajo contratos alternos. \\
Superposicion geometrica C &
18 & no clasifica target &
No aplica; parte de estados del candidato &
Mide persistencia geometrica acotada, no ocultedad. \\
\bottomrule
\end{longtable}

\section{Sistema y equilibrios}

Los analisis resumidos aqui usan el sistema de Chua no suave de Danca 2017:
\[
{}^C D_t^q x = \alpha(y-x-f(x)),\qquad
{}^C D_t^q y = x-y+z,\qquad
{}^C D_t^q z = -\beta y-\gamma z,
\]
con
\[
\alpha=8.4562,\quad \beta=12.0732,\quad \gamma=0.0052,\quad
m_0=-0.1768,\quad m_1=-1.1468,\quad q=\qmain.
\]
Los equilibrios calculados y usados por las pruebas son:
\[
\Eminus=(-6.588307887,-0.002836402,6.585471484),\quad
\Ezero=(0,0,0),\quad
\Eplus=(6.588307887,0.002836402,-6.585471484).
\]
El criterio de Matignon da \(\Ezero\) estable localmente y \(\Eminus,\Eplus\)
inestables para \(q=\qmain\). Por eso la verificacion de ocultedad debe
muestrear vecindades de los tres equilibrios, no solo observar separacion
visual del atractor.

\section{Atractor reportado por Danca}

\subsection{Contrato numerico}

La replicacion disponible usa el integrador Adams--Bashforth--Moulton para
Caputo con historia completa, sin truncamiento de memoria:
\[
h=0.05,\qquad t_{\mathrm{final}}=500,\qquad
t_{\mathrm{trans}}=250,\qquad \delta=0.01.
\]
El PDF usado no publica la condicion inicial de la Fig. 3; por ello se localizo
una semilla compatible, identificada como
\status{machado\_mu2\_seed\_project}. La semilla produjo una trayectoria
\status{bounded\_nontrivial} y fue usada como referencia numerica del atractor
de Danca.

\begin{table}[h]
\centering
\caption{Prueba local tipo Danca: trayectorias desde vecindades de los equilibrios inestables.}
\small
\begin{tabular}{lrrrrl}
\toprule
Equilibrio & Distancia inicial & Trayectorias & \target{} & Infinito & Decision \\
\midrule
\(\Eplus\)  & \(1.0\times10^{-2}\) & 100 & 54 & 46 & \status{not\_supported\_by\_Danca\_delta\_test} \\
\(\Eminus\) & \(1.0\times10^{-2}\) & 100 & 48 & 52 & \status{not\_supported\_by\_Danca\_delta\_test} \\
\midrule
Total & -- & 200 & 102 & 98 & no soportado como oculto \\
\bottomrule
\end{tabular}
\end{table}

Todas las condiciones iniciales de esta prueba estan a distancia
\(\delta=0.01\) del equilibrio correspondiente. Entre las trayectorias que
terminaron en \target{}, la distancia final al equilibrio mas cercano quedo
aproximadamente entre \(0.275\) y \(0.738\), con medianas \(0.667\) para
\(\Eplus\) y \(0.680\) para \(\Eminus\). Esto indica que la trayectoria sale de
una vecindad local muy pequena y termina lejos del equilibrio, en la region
clasificada como atractor candidato.

\subsection{Corte de cuenca Danca vs candidato del proyecto}

Tambien se ejecuto un corte \(xy\) de \(101\times101=10201\) condiciones
iniciales para comparar el contrato ABM de Danca contra el clasificador EFORK
del candidato del proyecto. Los cortes no son metodo-identicos: Danca usa ABM
con historia completa y el proyecto usa EFORK con memoria finita.

\begin{table}[h]
\centering
\caption{Conteos del corte de cuenca \(101\times101\).}
\small
\begin{tabular}{lrrrrr}
\toprule
Contrato & Equilibrio & Target + & Target - & Infinito & Desconocido \\
\midrule
Danca ABM & 1 & 514 & 462 & 8747 & 477 \\
Proyecto EFORK & 0 & 1 & 2 & 9097 & 1101 \\
\bottomrule
\end{tabular}
\end{table}

La lectura es que el corte muestra regiones que alcanzan clases tipo target,
pero el criterio de ocultedad queda decidido por las vecindades de equilibrio:
con \(\delta=0.01\), Danca no queda soportado como oculto en esta replicacion.

\section{Ruta Lur'e sesgada}

\subsection{Busqueda armonica}

La exploracion Lur'e sesgada se ejecuto con \(q=\qmain\). La funcion
descriptiva se uso como generador de semillas de primer armonico, no como
prueba de ciclos exactos en Caputo. La evaluacion principal produjo:

\begin{table}[h]
\centering
\caption{Resumen de la exploracion Lur'e sesgada.}
\small
\begin{tabular}{lr}
\toprule
Cantidad & Valor \\
\midrule
Evaluaciones crudas & 5208 \\
Candidatos filtrados & 121 \\
Semillas generadas & 968 \\
Filas de continuacion registradas & 38 \\
Supervivientes de continuacion & 4 \\
Candidatos que pasaron filtro temprano desde \(E_+\) & 0 \\
Supervivientes robustos registrados & 0 \\
\bottomrule
\end{tabular}
\end{table}

El mejor candidato por residual fue
\status{lure\_biased\_q\_0p99980\_rank\_0001}, con
\[
A=5.8517686,\quad \sigma_0\approx -9.75\times10^{-11},\quad
\omega=2.0402861,\quad |R|=6.82\times10^{-13},\quad
\rho_H=0.123874.
\]
El residual pequeno solo confirma consistencia de la semilla armonica de primer
orden. La validacion posterior debe mirar las trayectorias de Caputo.

\subsection{Continuacion y filtro temprano}

Cuatro candidatos sobrevivieron la continuacion como
\status{bounded\_nontrivial}, con memoria transportada y confiabilidad alta:
\status{rank\_0001}, \status{rank\_0003}, \status{rank\_0004} y
\status{rank\_0005}. Sin embargo, el filtro temprano desde equilibrios encontro
contactos \target{}.

\begin{table}[h]
\centering
\caption{Contactos desde vecindades de equilibrio en la ruta Lur'e.}
\small
\begin{tabular}{llrrl}
\toprule
Candidato & Equilibrio & Distancia inicial & Direcciones \target{} & Conclusion \\
\midrule
\status{rank\_0001} & \(\Eplus\) & \(1.0\times10^{-5}\) & 1 & bloquea ocultedad \\
\status{rank\_0004} & \(\Eplus\) & \(1.0\times10^{-5}\) & 8 & bloquea ocultedad \\
\status{rank\_0004} & \(\Eminus\) & \(1.0\times10^{-5}\) & 4 & bloquea ocultedad \\
\status{rank\_0004} & \(\Ezero\) & \(1.0\times10^{-5}\) & 0 & sin contacto target en ese equilibrio \\
\bottomrule
\end{tabular}
\end{table}

Para \status{rank\_0001}, basta una sola direccion desde \(\Eplus\), a
distancia \(10^{-5}\), para impedir una declaracion de atractor oculto. Para
\status{rank\_0004}, las direcciones \target{} desde \(\Eplus\) y \(\Eminus\)
son todavia mas fuertes.

\section{Ruta Machado}

\subsection{Barridos de semillas}

La ruta Machado usa una funcion descriptiva fraccionaria auxiliar como
generador de semillas. El parametro \(\mu\) de Machado no es el orden de
Caputo: \(q\) se mantuvo en \(\qmain\) para la corrida principal.

\begin{table}[h]
\centering
\caption{Barridos Machado disponibles.}
\small
\begin{tabularx}{\textwidth}{lrrrrX}
\toprule
Barrido & \(q\) & Registros & OK & Fallidos & Lectura \\
\midrule
\status{machado\_sweep} & 0.9998 & 552 & 264 & 288 &
Barrido principal. Todos los registros OK quedaron
\status{not\_supported\_by\_sample}; los mejores tuvieron al menos un contacto
\target{}. \\
\status{machado\_sweep\_fast} & 0.99 & 72 & 36 & 36 &
Barrido auxiliar/historico con otro \(q\). No debe mezclarse con las
conclusiones de \(q=0.9998\). \\
\bottomrule
\end{tabularx}
\end{table}

En el barrido principal, la regla de seleccion fue: menor numero de contactos
\target{}, desempate por mayor rango en \(x\) y menor norma de semilla. Los dos
mejores candidatos fueron:

\begin{table}[h]
\centering
\caption{Mejores candidatos Machado del barrido principal.}
\small
\begin{tabular}{lrrrrrl}
\toprule
Candidato & \(\mu\) & \(\theta\) & \(\omega\) & \(A\) & \target{} muestra & Estado inicial \\
\midrule
\status{branch\_0\_mu\_4p00000\_theta\_0p00000}
& 4.0 & 0.0 & 2.040286 & 1.714948 & 1 & \status{not\_supported\_by\_sample} \\
\status{branch\_0\_mu\_2p00000\_theta\_3p92699}
& 2.0 & 3.926991 & 2.040286 & 2.628434 & 1 & \status{not\_supported\_by\_sample} \\
\bottomrule
\end{tabular}
\end{table}

\subsection{Corrida 1 refinada}

La primera reverificacion fina intento reemplazar la lectura de barrido por
muestreos mas densos en vecindades. La corrida planificada era grande:
49344 simulaciones para la etapa A si se cubrian candidatos, radios, modos de
muestreo y equilibrios. En la ejecucion disponible se completaron 1622 filas
crudas, todas para el candidato \(\mu=4,\theta=0\) en vecindades de \(\Ezero\).

\begin{table}[h]
\centering
\caption{Estado de Corrida 1 refinada.}
\small
\begin{tabularx}{\textwidth}{L{0.34\textwidth}rrrrX}
\toprule
Candidato & Trayectorias & Radios efectivos & \target{} & Distancia inicial & Decision \\
\midrule
\status{branch\_0\_mu\_4p00000\_theta\_0p00000}
& 1622 & \(10^{-5},3\times10^{-5}\) & 0 &
bola/cascaron alrededor de \(\Ezero\) &
\status{compatible\_with\_hiddenness\_under\_tested\_radii}; no es \hiddenverified. \\
\status{branch\_0\_mu\_2p00000\_theta\_3p92699}
& 0 & -- & 0 & -- &
\status{not\_evaluated\_cost\_guard}; no se puede concluir desde esta corrida. \\
\bottomrule
\end{tabularx}
\end{table}

La compatibilidad de la primera fila era local y parcial: cubria solamente los
radios y trayectorias que llegaron a completarse. Por eso se ejecuto una ruta
dirigida posterior.

\subsection{Verificacion dirigida Machado con memoria corta}

La verificacion dirigida uso un contrato EFORK de memoria corta:
\[
q=0.9998,\qquad h=0.01,\qquad L_m=10,\qquad t_{\mathrm{final}}=3000
\]
para reconstruir los atractores de referencia. Se hicieron dos reconstrucciones
de atractor, una por candidato, ambas \status{bounded\_nontrivial}. Luego se
probaron direcciones dirigidas desde vecindades de equilibrio y se reprodujo el
contacto con un contrato mas estricto.

\begin{table}[h]
\centering
\caption{Prueba dirigida y reproduccion de contactos Machado.}
\small
\begin{tabularx}{\textwidth}{L{0.34\textwidth}lrrrrX}
\toprule
Candidato & Eq. & Distancia inicial & Tests dirigidos & \target{} & Reproducidos & Resultado \\
\midrule
\status{branch\_0\_mu\_4p00000\_theta\_0p00000}
& \(\Eminus\) & \(1.0\times10^{-5}\) & 1 & 1 & 1 &
\status{not\_supported\_by\_equilibrium\_neighborhood\_test}. \\
\status{branch\_0\_mu\_2p00000\_theta\_3p92699}
& \(\Eminus\) & \(1.0\times10^{-5}\) & 1 & 1 & 1 &
\status{not\_supported\_by\_equilibrium\_neighborhood\_test}. \\
\bottomrule
\end{tabularx}
\end{table}

La direccion fue \status{axis\_p\_x} desde \(\Eminus\). En ambos casos la
condicion inicial estaba a distancia \(9.9999999996\times10^{-6}\) de
\(\Eminus\). La trayectoria \target{} dirigida termino a distancia aproximada
2.795 del equilibrio mas cercano, es decir, no se quedo localmente pegada al
equilibrio; salio de la vecindad y alcanzo la geometria del candidato. La
reproduccion B uso \(h=0.005\), \(L_m=1\), \(t_{\mathrm{final}}=3000\), y
volvio a dar \target{} para ambos candidatos. Por eso la robustez larga se
omitio: ya habia un contacto reproducido desde equilibrio.

\section{Prueba top-3 con cascarones esfericos}

Para los tres candidatos que se estaban comparando
\[
\status{branch\_0\_mu\_4p00000\_theta\_0p00000},\quad
\status{branch\_0\_mu\_2p00000\_theta\_3p92699},\quad
\status{lure\_biased\_q\_0p99980\_rank\_0001},
\]
se lanzaron cascarones esfericos alrededor de \(\Ezero,\Eplus,\Eminus\) con
radios
\[
10^{-5},\quad 3\times10^{-5},\quad 10^{-4},\quad 3\times10^{-4},\quad 10^{-3}.
\]
Cada combinacion candidato--equilibrio--radio tuvo 100 condiciones iniciales,
por lo que cada candidato uso \(3\times5\times100=1500\) trayectorias y el
experimento completo tuvo 4500 trayectorias.

\begin{table}[h]
\centering
\caption{Resultado agregado de los cascarones esfericos top-3.}
\small
\begin{tabular}{lrrrrr}
\toprule
Candidato & Trayectorias & \target{} & Equilibrio & Infinito & Desconocido \\
\midrule
\status{branch\_0\_mu\_4p00000\_theta\_0p00000} & 1500 & 0 & 275 & 498 & 727 \\
\status{branch\_0\_mu\_2p00000\_theta\_3p92699} & 1500 & 0 & 277 & 512 & 711 \\
\status{lure\_biased\_q\_0p99980\_rank\_0001} & 1500 & 0 & 284 & 502 & 714 \\
\midrule
Total & 4500 & 0 & 836 & 1512 & 2152 \\
\bottomrule
\end{tabular}
\end{table}

Esta prueba es compatible con ocultedad bajo esos cascarones discretos, pero no
domina a las pruebas dirigidas: para los candidatos Machado ya existe un
contacto reproducido desde \(\Eminus\) a distancia \(10^{-5}\), y para Lur'e
existe un contacto temprano desde \(\Eplus\) a distancia \(10^{-5}\). La
ausencia de \target{} en los cascarones aleatorios debe leerse como ``no se
observo contacto en esa muestra'', no como refutacion de los contactos
dirigidos.

Los archivos de geometria de esferas usados para graficar no agregan nuevas
clasificaciones. Muestrean, por visualizacion, hasta 36 condiciones iniciales
por combinacion dentro de la esfera de radio maximo \(10^{-3}\) y recortan
segmentos cortos de trayectoria. Esas figuras quedan pendientes de inclusion.

\section{Pruebas de robustez y superposicion}

Se ejecutaron dos clases de prueba de robustez para los tres candidatos top-3.
La primera intenta reclasificar el candidato bajo cambios numericos; la segunda
superpone nubes de trayectoria para medir persistencia geometrica sin decidir
ocultedad.

\begin{table}[h]
\centering
\caption{Robustez por reclasificacion EFORK en top-3.}
\small
\begin{tabular}{lrrrl}
\toprule
Candidato & Casos & Casos target & No target & Estado \\
\midrule
\status{branch\_0\_mu\_4p00000\_theta\_0p00000} & 4 & 0 & 4 & \status{not\_robust\_under\_tested\_contracts} \\
\status{branch\_0\_mu\_2p00000\_theta\_3p92699} & 4 & 0 & 4 & \status{not\_robust\_under\_tested\_contracts} \\
\status{lure\_biased\_q\_0p99980\_rank\_0001} & 4 & 0 & 4 & \status{not\_robust\_under\_tested\_contracts} \\
\bottomrule
\end{tabular}
\end{table}

La superposicion geometrica C lanzo 18 trayectorias: 3 candidatos por 6
contratos numericos. Los contratos incluyeron \(h=0.005,0.01,0.02\),
\(L_m=5,10,20\) y \(t_{\mathrm{final}}=1500,3000\). Esta prueba no clasifica
\target{} ni decide ocultedad. En ella todas las trayectorias fueron acotadas,
sin divergencia, sin colapso a equilibrio y con varianza no colapsada. El caso
\(L_m=20\) debe leerse como sensibilidad numerica, no como contrato STM32 de
memoria corta.

\section{Conclusiones operativas}

\begin{enumerate}
\item El atractor reportado por Danca fue reproducido como dinamica acotada no
trivial bajo ABM con historia completa, pero la prueba local \(\delta=0.01\)
produjo 102 contactos \target{} de 200 trayectorias desde \(\Eplus\) y
\(\Eminus\). En esta replicacion no queda soportado como oculto.

\item El mejor candidato Lur'e,
\status{lure\_biased\_q\_0p99980\_rank\_0001}, tiene residual armonico muy
pequeno y sobrevive continuacion EFORK, pero una condicion inicial a distancia
\(10^{-5}\) de \(\Eplus\) llega a \target{}. Por tanto no debe etiquetarse
como \hiddenverified.

\item Los dos mejores candidatos Machado de \(q=\qmain\) sobreviven como
atractores acotados no triviales, pero la verificacion dirigida encontro y
reprodujo contactos \target{} a distancia \(10^{-5}\) de \(\Eminus\). La
conclusion recomendada es documentarlos como candidatos autoexcitados o no
ocultos bajo el contrato probado, no como atractores ocultos.

\item La prueba de cascarones esfericos top-3 lanzo 4500 trayectorias y no
encontro \target{}. Esa evidencia es util para mapear geometria de cuencas, pero
no cancela los contactos dirigidos, que son mas especificos y reproducidos.

\item Las pruebas de robustez geometrica muestran persistencia de trayectorias
acotadas bajo cambios de \(h\), \(L_m\) y tiempo total, pero robustez del
atractor y ocultedad son criterios distintos. La robustez no convierte un
candidato con contacto desde equilibrio en oculto.
\end{enumerate}

\section{Artefactos consultados}

\begin{longtable}{@{}L{0.42\textwidth}L{0.50\textwidth}@{}}
\toprule
Archivo & Uso en este reporte \\
\midrule
\endhead
\codepath{outputs/danca2017_chua_abm_20260515_182354/danca_hiddenness_summary.csv} &
Conteos de trayectorias locales desde \(\Eplus\) y \(\Eminus\). \\
\codepath{outputs/danca2017_chua_abm_20260515_182354/danca_hiddenness_raw.csv} &
Distancias iniciales y clases finales de la prueba Danca. \\
\codepath{outputs/danca2017_chua_abm_20260515_182354/danca_reference_summary.json} &
Contrato ABM y semilla de referencia localizada. \\
\codepath{outputs/basin_compare_danca_project_h001_grid101_20260517/basin_comparison_summary.json} &
Conteos del corte de cuenca Danca/proyecto. \\
\codepath{outputs/lure_biased_multiparam_q09998/lure_biased_multiparam_summary.json} &
Resumen global de busqueda Lur'e. \\
\codepath{outputs/lure_biased_multiparam_q09998/continuation_survivors.csv} &
Supervivientes de continuacion Lur'e. \\
\codepath{outputs/lure_biased_multiparam_q09998/early_equilibrium_filter_summary.csv} &
Contactos tempranos desde vecindades de equilibrio en Lur'e. \\
\codepath{runs_machado_sweep_fast/chua_piecewise/machado_sweep/machado_sweep_summary.json} &
Barrido Machado principal con \(q=0.9998\). \\
\codepath{outputs/extended_search/corrida1/refined_hiddenness_decision.csv} &
Estado de Corrida 1 refinada. \\
\codepath{outputs/extended_search/machado_targeted_verification_lm10_20260515_182252/machado_targeted_decision.csv} &
Decision final de verificacion dirigida Machado. \\
\codepath{outputs/top3_machado_lure_sphere_C_20260517_v2/sphere_decision.csv} &
Decision de cascarones esfericos top-3. \\
\codepath{outputs/top3_machado_lure_sphere_C_20260517_v2/sphere_raw.csv} &
Conteos agregados por clase y distancias de radios esfericos. \\
\codepath{outputs/robustness_overlay_c_trajectories_20260517/robustness_overlay_metrics.csv} &
Persistencia geometrica bajo cambios de contrato numerico. \\
\bottomrule
\end{longtable}

\end{document}
